\section{Residual-adapted finite-panel compression}
\label{sec:panel-proof}

This appendix proves the Logarithmic row of Theorem~\ref{thm:compression}.
It uses the same selected-coordinate optimizer as Harmonic compression,
but constructs its source spaces from finitely many response curves.
The new argument concerns the complex-time domain of those curves.
The coordinate-selection, fitting, and cancellation arguments are reused
only after their hypotheses have been verified below.  All references in
this appendix are to proofs included in this paper.

\subsection{Contract, exact coefficients, and width qualification}
\label{sec:panel-contract}

Fix the training inputs and \(p\) additional inputs before initialization.
Only the first \(m\) inputs have labels or enter training sums.  Write
\(v_i=x_i/\sqrt d\), so \(\|v_i\|_2=1\), and use the canonical
Gaussian initialization, zero readout, mobilities, and physical time of
the main text.  The passive labels are neither required nor used.  The
activation assumptions are those of
\S\ref{int:source-foundations}: holomorphy on \(|\operatorname{Im}z|<a\),
reality on the real axis, and bounded first derivative on that strip.
Activation values need not be bounded.  Put
\begin{align*}
 b&=\max_j|\phi_j(0)|,&
 s&=\max\{1,\sup_{j,|\operatorname{Im}z|\le a/2}|\phi_j'(z)|\},\\
 t_2&=\max\{1,\sup_{j,|\operatorname{Im}z|\le a/2}|\phi_j''(z)|\},&
 \beta&=\max(10,1+b,s,t_2,16/a).
\end{align*}
The second derivative bound is finite by the strip hypothesis and
Cauchy's formula.  Throughout this appendix the local abbreviations are
\[
 \lambda=\gamma/m,\qquad Y=\|y\|_2/\sqrt m>0,\qquad
 S=16Y/\lambda,\qquad \ell=\log(en),\qquad
 \nu=\lambda/4,\qquad T_0=32\ell/\lambda.
\]
The case \(Y=0\) has the exact stationary zero predictor and needs no
division by \(Y\).  The full label allowance is
\begin{equation}
 \frac{Y}{\lambda}\le
 \min\left\{\frac1{8H_d\sqrt{F_d}},
             \frac1{16H_c\sqrt{F_c}},\frac{S_*^{\rm src}}{16}\right\}.
 \label{eq:panel-labels}
\end{equation}
Here \(H_d,F_d,H_c,F_c\) are exactly the finite recurrences in
\eqref{int:eq-0054-harmonic-dense-coefficients}--\eqref{int:eq-0055-harmonic-fitting-coefficients},
and \(S_*^{\rm src}\) is the explicit minimum
\eqref{int:eq-0184-s-10}.  Thus \eqref{eq:panel-labels} is precisely
the existing Harmonic allowance, not a new smaller cap.  In particular
\(S\le1\).  The additional Legendre allowance is imposed only when
all three methods are asserted together.

For an explicit finite prescription we recall the source coefficients.
All symbols in the following recurrence are fixed scalars, not model
coordinates or approximation orders:
\begin{align*}
 H_1&=b+20s,&H_j&=b+10sH_{j-1},&
 k_j&=H_L(10s)^{L-j},&\tau_j&=sk_j,\\
 P_1&=3,&P_j&=H_{j-1}+10sP_{j-1}+1,&f_j&=sP_j,\\
 g&=\tau_1+\sum_{j=2}^L\tau_jH_{j-1},&
 r_j^{\rm src}&=P_jg,&q_j^{\rm src}&=f_jg,\\
 A_*&=2sP_L+2s^2\sum_{j=2}^Lk_jP_{j-1},&
 H_*&=A_*+t_2\sum_{j=1}^LP_j^2k_j,\\
 E_*&=t_2\sum_{j=1}^L(10s)^{2(j-1)}k_j.
\end{align*}
Continue with
\begin{align*}
 D_0&=(1+b+s)\{1+2H_*^2+4A_*^2(e^2-1)+E_*+s^2\max_jk_j^2\},\\
 K_{\rm src}&=128(D_0+1)\{1+32\max(1,\max_jH_j,\max_j\tau_j)\},\\
 e_1&=t_2r_1^{\rm src}P_1,\\
 e_j&=t_2r_j^{\rm src}P_j+
 s(q_{j-1}^{\rm src}+\tau_jH_{j-1}f_{j-1}+10e_{j-1}),\\
 T_Q&=8\max_jf_j\max_j(f_jH_*+e_j),\\
 U=U_{\rm fin}(S)&=\max\left\{4sK_{\rm src},\right.\\[-2mm]
 &\quad\left.\max_{j\ge2}2\left[sK_{\rm src}
 (H_{j-1}^2+f_{j-1}^2+ST_Q+S^2H_{j-1}q_{j-1}^{\rm src})
 +64q_{j-1}^{\rm src}+1\right]\right\}.
\end{align*}
These are the coefficients of \S\ref{int:source-explicit-recurrences}
and \S\ref{int:dense-finite-query-source}; the superscripts on
\(r_j^{\rm src},q_j^{\rm src}\) prevent confusion with residuals and
selected widths.  Define
\begin{equation}
 \begin{split}
 \mathcal K&=H_L^2+S^2\left[\tau_1^2+
                  \sum_{j=2}^L\tau_j^2H_{j-1}^2\right],\qquad
 M_0=10\max_j(H_j,\tau_j),\\
 r_n&=\frac{a}{64YSU\sqrt\ell},\qquad
 D_W=\max\{\tau_1,\max_{j\ge2}\tau_jH_{j-1}\}.
 \end{split}
 \label{eq:panel-source-coefficients}
\end{equation}
In the source budget \eqref{int:eq-0184-s-10}, denote its exponential
coefficient by \(\eta_{\rm src}\) and its ceiling by
\(\mathcal B=1024e^2L\), to distinguish them from an approximation
tolerance.  Both retain their exact definitions there.

The theorem is eventual for each fixed admissible problem and each
fixed confidence.  In addition to the unchanged initialization and
source-event onsets, the numerical gates used below are
\begin{gather}
 n^{-1}\le\min(1,Y,S),\qquad
 \ell\ge\max(e^2,2\mathcal B),\qquad
 \sqrt\ell\ge \frac{a}{64YSU}
       \max\{8,\lambda,4\mathcal K/\log2,32YSD_W\},
 \label{eq:panel-old-gates}\\
 \max\left\{4\mathcal K r_n,
     \frac{2C_g}{\sqrt{\mathcal K}}B_nYr_n\right\}<\log(3/2),
 \qquad B_n=1+\frac{S^2}{\eta_{\rm src}}\log(e+\ell),
 \label{eq:panel-new-gate}
\end{gather}
where the finite coefficient \(C_g\) is given in
\eqref{eq:panel-gradient-derivative} below.  The fitted-tail and degree
gates are \eqref{eq:panel-tolerance} and \eqref{eq:panel-degree-gate}.
Every displayed gate holds eventually at fixed parameters.  The
probability onset is qualitative: no effective polynomial threshold,
uniform growing-panel statement, or simultaneous event over infinitely
many independent widths is asserted.

\subsection{The enlarged analytic domain}
\label{sec:panel-domain}

The independently proved real fitting theorem gives
\(\rho(t)=\|r(t)\|_2/\sqrt m\le Ye^{-\nu t}\), physical operator
caps strictly below nine, and \(2\int_0^\infty\rho\le S/2\).
These facts also hold for each fixed-size rectangular cavity on its
own initialization event.  This last statement uses the zero-embedding
calculation \eqref{int:eq-0286-ic-71}--\eqref{int:eq-0288-ic-73}; it
does not identify a rectangular cavity covariance with the full one.

Put
\begin{equation}
 h(t)=\frac1{2\mathcal K}\log(1+2\mathcal K r_ne^{\nu t}),\quad
 h_0=h(0),\quad d_h=\frac\nu{2\mathcal K},\quad
 T_+=T_0+h(T_0)+r_n.
 \label{eq:panel-flaring-height}
\end{equation}
Since \(\lambda\le H_L^2\le\mathcal K\),
\(0\le h'\le d_h\le1/8\), \(h_0\le r_n\), and
\(h(t)\le r_n+d_ht\).  The inequality \(\lambda\le H_L^2\)
follows by the population RMS recurrence and
\(\gamma\le Q^{(L)}_{aa}\).  For \(0\le s_0\le1\) use
\[
 D_{s_0}=\{x+iy:-s_0h_0\le x\le s_0T_+,\quad
                     |y|\le s_0h(\max(x,0))\}.
\]
These nested, vertically fibered domains are not assumed convex.
Each point is reached by a real segment from zero and then a vertical
segment.  For negative real part the combined exceptional length is
at most \(2h_0\).

Each cavity is stopped independently at its own first-exit level
\(s_c\).  Use the source stops of
\S\ref{int:source-local-insertion}: full and cavity exponential budgets
\(\mathcal B,2\mathcal B\), provisional maxima and response caps
twice and four times their claimed values, pole margins \(3a/8,7a/16\),
and the unchanged operator, activity, row-size and passive-preactivation
stops.  A failed own initialization gives identically zero reference
coefficients.  This convention precedes integration over omitted roots.

Here is the globally defined extension needed for conditional Gaussian
arguments.  For real-interval clipping write
\[
 X_{s_0}(x)=\operatorname{clip}_{[-s_0h_0,s_0T_+]}(x),\qquad
 P_{s_0}(x+iy)=X_{s_0}(x)+i\operatorname{clip}_{[-s_0h(X_{s_0}(x)_+),
                  s_0h(X_{s_0}(x)_+)]}(y).
\]
This is a retraction onto \(D_{s_0}\), measurable in \(s_0\).
With \(a_0(z)=\operatorname{clip}_{[0,T_+]}(\operatorname{Re}z)\),
direct clipping gives
\begin{equation}
 P_{s_0}a_0(z)=a_0(P_{s_0}z),\qquad
 \operatorname{Lip}(P_{s_0})\le1+d_h.
 \label{eq:panel-clipping}
\end{equation}
For the Lipschitz assertion the coordinate increments satisfy
\(|\Delta X|\le|\Delta x|\) and
\(|\Delta\operatorname{Im}P|\le|\Delta y|+d_h|\Delta x|\);
the resulting shear has norm at most \(1+d_h\).
For a coefficient with derivative bound on \(D_{s_0}\), integrate
along the image under \(P_{s_0}\) of a straight segment.  This proves
the corresponding global Lipschitz bound without convexity.
Extend each stopped coefficient as \(b_c(P_{s_c}z)\) on the deterministic
rectangle
\(E_n=[-h_0,T_+]+i[-h(T_+),h(T_+)]\).
Its real anchors traverse a monotone real prefix and then freeze.
Equation~\eqref{eq:panel-clipping} supplies the same commuting-anchor
identity as \eqref{int:eq-0122-ic-stop}.

Local holomorphic ODE continuation and uniqueness apply on every strict
prefix: finite parameter bounds and a strict activation-strip margin
permit extension, and overlapping real-anchor/vertical patches agree.
The domain retracts onto its real interval and is simply connected.
At a first exit, the strict margin below the half-strip permits one-sided
derivative estimates.  A strict improvement of all stops therefore
extends the solution across the compact boundary.  The clipped extensions
themselves need not be holomorphic outside their own stopped domains.

\subsection{The derivative estimate and the noncircular propagator bound}
\label{sec:panel-propagator}

Work first under each cavity's provisional response stop
\(\max|R_{ba,i}^{(j)}|\le4SU\sqrt\ell\), where
\(R_{ba}^{(j)}=D_\Theta z_b^{(j)}\nabla_\Theta F_a\),
\(F_a=nf_n(v_a)\), and
\(\Theta=(A,\sqrt nW^{(2)},\ldots,\sqrt nW^{(L)},w)\).
The exact derivative equations give
\[
 \frac{\|\dot z_a^{(j)}\|_2}{\sqrt n}\le2\rho S r_j^{\rm src},
 \qquad \|\dot z_a^{(j)}\|_\infty\le8\rho SU\sqrt\ell.
\]
Set \(N_*=\max_j\{r_j^{\rm src},4U\}\).
Interpolation and the stopped exponential carrier budget imply, for
\(u\ge4\),
\[
 \frac{\|\dot z_a^{(j)}\|_{2u/(u-2)}}{n^{(u-2)/(2u)}}
       \le2\rho SN_*\ell^{1/u},\qquad
 \frac{\|k_a^{(j)}\|_u}{n^{1/u}}
       \le\frac S{\eta_{\rm src}}u(2\mathcal B)^{1/u}.
\]
H\"older consequently bounds the normalized 2-norm of
\(k_a^{(j)}\odot\dot z_a^{(j)}\) by
\(2\rho S^2N_*u(2\mathcal B\ell)^{1/u}/\eta_{\rm src}\).
Choose \(u=\max(4,\log(2\mathcal B\ell))\).  Under
\eqref{eq:panel-old-gates}, the exponential factor is at most \(e\)
and \(u\le2\log(e+\ell)\).  Differentiation of the backward
recursion then gives
\begin{gather*}
 J_L=2sH_L+4et_2N_*,\qquad
 J_j=10sJ_{j+1}+2s^3k_{j+1}^2H_j+4et_2N_*,\\
 J_* =\max_jJ_j,\qquad
 \frac{\|\dot h_a^{(j)}\|_2}{\sqrt n}\le2\rho Sq_j^{\rm src},
 \qquad
 \frac{\|\dot\delta_a^{(j)}\|_2}{\sqrt n}\le J_*B_n\rho.
\end{gather*}
The four terms are respectively the changed readout, propagated upper
derivative, changed mixer, and changed activation gate.  This derivation
uses only the provisional stops, not their later Gaussian improvement.

For the normalized gradient matrix
\(G=(mn)^{-1/2}[\nabla_\Theta F_1,\ldots,\nabla_\Theta F_m]\),
its column blocks are \(m^{-1/2}\) times
\(\delta_a^{(1)}v_a^\top/\sqrt n\),
\(\delta_a^{(j)}h_a^{(j-1)\top}/n\), and
\(h_a^{(L)}/\sqrt n\).  Differentiating these products proves
\begin{equation}
 \begin{split}
 \|G\|_{\rm op}&\le\sqrt{\mathcal K},\qquad
 \|\dot G\|_{\rm op}\le C_gB_n\rho,\\
 C_g&=J_*\left(1+\sum_{j=2}^LH_{j-1}\right)
       +2S^2\sum_{j=2}^L\tau_jq_{j-1}^{\rm src}+2Sq_L^{\rm src}.
 \end{split}
 \label{eq:panel-gradient-derivative}
\end{equation}
Bounding the Frobenius norm after division of each column by \(\sqrt m\)
introduces no sample factor.

The exact residual equation is \(\dot r=-2G^\top Gr\), and the
residual-free variational base is \(\dot P=-2GG^\top P\).
Transposes here are algebraic.  On a vertical segment from a real
anchor \(t\ge0\), norm differentiation yields
\begin{equation}
 \rho(t\pm iv)\le Ye^{-\nu t}e^{2\mathcal K v},\qquad
 \int_0^{h(t)}\rho(t\pm iv)\,dv\le Yr_n.
 \label{eq:panel-vertical-activity}
\end{equation}
The integral follows by substituting
\(e^{2\mathcal K h(t)}=1+2\mathcal K r_ne^{\nu t}\).
It also implies \(\rho\le Y+2\mathcal K Yr_n\le2Y\).

At the real anchor \(A_t=2G(t)G(t)^\top\) is real symmetric;
therefore \(-iA_t\) generates unitary motion.  Freeze this generator
along the vertical segment.  By \eqref{eq:panel-gradient-derivative},
\[
 \|2G(t\pm iv)G(t\pm iv)^\top-A_t\|_{\rm op}
 \le4\sqrt{\mathcal K}C_gB_n\int_0^v\rho(t\pm iu)\,du.
\]
Conjugating by the frozen unitary gives a propagator norm at most
\begin{equation}
 \exp\left(4\sqrt{\mathcal K}C_gB_n
      \int_0^{h(t)}(h(t)-v)\rho(t\pm iv)\,dv\right)
 \le\exp\left(\frac{2C_g}{\sqrt{\mathcal K}}B_nYr_n\right).
 \label{eq:panel-unitary-bound}
\end{equation}
Indeed the double integral is at most \(Yr_n/(2\mathcal K)\)
by integrating the exponential in \eqref{eq:panel-vertical-activity}.
Products on ordered disjoint subsegments have the same bound: use the
same frozen generator and add their nonnegative growth integrals.
Forward real segments contract.  The negative-real corner instead has
length at most \(2r_n\) and costs at most \(e^{4\mathcal K r_n}\).
Thus \eqref{eq:panel-new-gate} supplies a product bound below \(3/2\),
strictly within the factor two required in
\eqref{int:eq-0126-ic-8} and \eqref{int:eq-0169-ic-51}.
This is established \emph{before} Gaussian insertion.

The extra activity is at most \(8Yr_n\).  The original gates give
\(8Yr_n\le S/2\) and \(8YSD_Wr_n\le1/4\), preserving activity
cap \(S\), strict operator margins inside ten, and
\(\rho/S\le\lambda/8\).  These are all the physical hypotheses
of \eqref{int:eq-0121-ic-4}.  Since
\(B_nr_n=O(\log(e+\ell)/\sqrt\ell)\), the added propagator gate
is eventual without a new label restriction.

\subsection{Transfer of the source proof to the larger domain}
\label{sec:panel-source-transfer}

We now verify the domain-dependent hypotheses of the complete insertion
proof \S\ref{int:source-local-insertion}.  Its parameter derivative,
nonlinear-remainder, and Gaussian-form identities are unchanged.
First, canonical anchor contours have length at most
\((64/\lambda+4r_n\sqrt\ell)\ell\).  Replacing the contour-length
coefficient in \eqref{int:eq-0128-ic-10} by this fixed number leaves
the derivative recurrences \eqref{int:eq-0123-ic-5}--\eqref{int:eq-0135-ic-17}
and the passive response recurrences
\eqref{int:eq-0154-ic-36}--\eqref{int:eq-0161-ic-43} of size
\(n^{1/4000}\operatorname{poly}(\ell)\), with an additional
\(\sqrt n\) for raw terminal derivatives.  Their inputs are precisely
the unchanged operator/RMS/budget bounds, residual activity, provisional
maxima, and base product bound just verified.

Second, the one-dimensional control nets of
\eqref{int:eq-0162-ic-44} still have logarithmic cardinality
\(n^{5/8}\operatorname{poly}(\ell)\): control amplitudes and speeds
remain respectively polynomial-logarithmic and
\(\sqrt n\operatorname{poly}(\ell)\) along those contours.
The globally clipped coefficient paths have the same moduli up to
\(1+d_h\le9/8\).  A deterministic endpoint mesh in \(E_n\), whose
sides are \(O(\ell)\), has polynomial size.  The finite panel and
any fixed deletion-count unions therefore cost at most \(n^{0.65}\)
in logarithmic cardinality eventually, whereas the conditional Gaussian
linear/quadratic tails in \eqref{int:eq-0163-ic-45} have exponent
\(n^{0.78}\).  The same uniform event consequently fails with
probability at most \(e^{-n^{0.7}}\) eventually.  Deterministic
one-dimensional control histories are netted before substituting actual
analytic histories; no arbitrary two-dimensional control extension is
being assumed.

Third, the remainder identities
\eqref{int:eq-0140-ic-22}--\eqref{int:eq-0161-ic-43} use joining
segments inside the same safe activation half-strip.  Their Duhamel
bound \eqref{int:eq-0164-ic-46} uses the same activity and propagator
bound.  With the original cutoffs \(N=n^{1/100}\),
\(d_0=n^{-1/10}\), \(u_0=n^{-1/25}\), its remainder remains
\[
 n^{1/4000}\operatorname{poly}(\ell)
 [n^{-8/100}+u_0^2+n^{-48/100}+n^{-1/2}]=o(u_0).
\]
The common-prefix comparisons are therefore still coordinate discrepancy
\(n^{-1/30}\) and whole-vector discrepancy \(2n^{1/100}\).
The direct passive reverse observable in
\eqref{int:eq-0157-ic-39} is included; it is not replaced by an
independent reverse operator.

Fourth, the trace expansion
\eqref{int:eq-0166-ic-48}--\eqref{int:eq-0175-ic-57} requires products
over ordered disjoint base-propagator segments, exactly as verified in
\eqref{eq:panel-unitary-bound}.  Its ordered residual integrals remain
bounded by \(S^h/h!\).  Thus its Schatten coefficients, direct exterior
trace, learned row/column terms, and sample-RMS factors are unchanged.
In particular the scalar absorption \eqref{int:eq-0190-s-16} and
the coefficient \(U_{\rm fin}\) above are unchanged.

The scalar Gaussian maximum grid now has at most
\(O(n^4\ell^2)\le n^5\) points eventually.  Fixed panel, layer,
driver, and root indices fit the original enclosing count \(n^{10}\).
The tail union \(4n^{10}e^{-1024\ell}\) and off-grid error
\(n^{-3/2}\operatorname{poly}(\ell)\) vanish.  Hence the full
provisional caps improve to
\begin{equation}
 \max|z_{a,i}^{(j)}|\le K_{\rm src}\sqrt\ell,\qquad
 \max|k_{a,i}^{(j)}|\le SK_{\rm src}\sqrt\ell,\qquad
 \max|R_{ba,i}^{(j)}|\le SU\sqrt\ell.
 \label{eq:panel-improved-caps}
\end{equation}
The last estimate has each training driver and each declared evaluated
input.  Passive backward budgets are unnecessary.  The passive
preactivation stop improves by its initial finite-panel Gaussian maximum
and its activity integral.  Combining
\eqref{eq:panel-vertical-activity} and \eqref{eq:panel-improved-caps}
bounds displacement from a real anchor by \(a/32\); the negative corner
gives \(a/8\).  The pole stop improves strictly.

A cavity cannot stop first on the successful full prefix: coordinate
comparison transfers its maxima and pole margin, while its budget is at
most \(e^{2\eta_{\rm src}n^{-1/30}}\mathcal H_a+O_u(1/n)
<2\mathcal B\) for fixed deletion count \(u\).  Operator and activity
margins were already proved deterministically.  This verifies the
strict independent-cavity transfer of \eqref{int:eq-0122-ic-stop}.
No Gaussian law has been conditioned on full-network survival.

It remains to remove the budgets on complete independently stopped
paths, not just that prefix.  For a forward coefficient normalized by
\(\sqrt n\), its complex-minus-real radius is at most
\(2S\max_jq_j^{\rm src}Yr_n\).  For a backward coefficient normalized
by \(S\sqrt n\), it is at most \(J_*B_nYr_n/S\).  The negative
corner costs at most four times these amounts.  Commutation
\eqref{eq:panel-clipping} thus gives the complete-path radius
\[
 D_n=4Yr_n\max\{2S\max_jq_j^{\rm src},J_*B_n/S\}
      =O(\log(e+\ell)/\sqrt\ell).
\]
Its global normalized modulus is \(O(B_n)\).  Applying the proved
Gaussian rectangle estimate \eqref{int:eq-0292-ic-77} gives mean
\(O(\log(e+\ell)^{3/2}/\sqrt\ell)\), tail radius \(2D_n\),
and every fixed exponential moment tending to one.  Its sample-RMS
version requires no independence among samples.

For singleton/common-cavity differences, extend each path with its own
clipping map, subtract, and project onto the deterministic ball of radius
\(4n^{1/100}\).  Both paths omit the root paired with that difference;
projection preserves this independence and is the identity on the
successful prefix.  The normalized Gaussian radius is
\(4n^{-49/100}\), its modulus is polynomial-logarithmic, and
\eqref{int:eq-0292-ic-77} gives fixed exponential moments tending to one,
as in \eqref{int:eq-0293-ic-78}.

The real reference paths retain their original activity modulus and
squared-exponential estimate \eqref{int:eq-0194-s-20}.  Distinct omitted
roots are conditionally independent given the common cavity.  Separate
their main-reference products from all correction products by
Cauchy--Schwarz and then H\"older at each fixed moment order.  The
one-neuron main moment and scalar absorption retain the limiting base
sixteen of \eqref{int:eq-0203-s-29}.  Joint independence of correction
paths is neither used nor asserted.  Indicators of the full event are
dropped only from nonnegative, globally defined cavity bounds.
For a tuple of moment order \(u\) with only \(k<u\) distinct roots,
the normalized multiplicity is \(O_u(n^{k-u})\).  Repeated weights
cost at most \(e^{2u\eta_{\rm src}K_{\rm src}\sqrt\ell}=n^{o(1)}\)
by \eqref{eq:panel-improved-caps}, so their contribution vanishes.
The local exceptional events have superpolynomially small probability.
Markov's inequality and the fixed sample/layer union now give
\[
 \limsup_{n\to\infty}\Pr\{\text{a full budget is hit}\}
 \le mL(16L/\mathcal B)^u.
\]
Take the width limit at each fixed positive integer \(u\), then the
infimum over \(u\).  Since \(16L/\mathcal B<1\), the probability
tends to zero.  All other stops have strict improvements.  We have
therefore proved holomorphy on a neighborhood of \(D_1\), with the
original coordinate bound \(M_0\sqrt n\), for all required finite-panel
sources, with probability tending to one.  This proves the new source
event, rather than assuming it from the smaller rectangular event.

\subsection{Adaptive polynomial sources and initialization-only coefficients}
\label{sec:panel-compiler}

At each layer include the curves
\[
 h^{(j)}(t,v_i),\ W_0^{(j)}h^{(j-1)}(t,v_i)\quad(i\le m+p),
 \qquad
 \delta_a^{(j)}(t),\ W_0^{(j+1)\top}\delta_a^{(j+1)}(t)
       \quad(a\le m),
\]
omitting nonexistent boundary images.  There are at most \(2(2m+p)\)
curves per layer.  Linear activation growth
\(|\phi_j(z)|\le b+s|z|\), the source RMS bounds, and initialized
operator norm eight imply their coordinate bound \(M_0\sqrt n\).

Set \(\varrho(t)=h(t)/[2(1+d_h)]\) for \(0\le t\le T_0\).
Its closed disk about \(t\) lies in \(D_1\): for nonnegative real
coordinate \(x\), \(h(x)\ge h(t)-d_h\varrho(t)\ge\varrho(t)\).
If the disk reaches a negative coordinate, then \(t<\varrho(t)\)
and \(h(t)\le h_0+d_ht\) give \(\varrho(t)<h_0\).
The right endpoint is below \(T_0+h(T_0)<T_+\).
The radius is positive, nondecreasing, and Lipschitz with constant
\(d_h/[2(1+d_h)]\).

Starting at \(t_0=0\), set
\(t_{b+1}=\min(T_0,t_b+\varrho(t_b)/2)\) until \(t_J=T_0\).
For any nondecreasing radius with Lipschitz constant \(C_r\), every
full panel contributes at least \(1/(2+C_r)\) to
\(\int dt/\varrho(t)\); hence
\(J\le1+(2+C_r)\int_0^{T_0}dt/\varrho(t)\).
For our radius this yields
\begin{equation}
 J\le1+40960\frac Ua\left(\frac Y\lambda\right)^2\sqrt\ell
       +80\frac{\mathcal K}{\lambda}
                  \log\left(1+\frac{4\ell}{\log2}\right).
 \label{eq:panel-panel-count}
\end{equation}
For completeness, it suffices to bound \(J\) by
\(1+5\int_0^{T_0}dt/h(t)\).  Split where
\(2\mathcal K r_ne^{\nu t}=1\).  Before that point use
\(\log(1+x)\ge x/2\); afterward its logarithm grows at least
with slope \(\nu/2\).  Thus the integral is at most
\(2/(\nu r_n)+(4\mathcal K/\nu)
\log(1+\nu T_0/(2\log2))\), giving the displayed constants.

Given source tolerance \(0<\eta\le\min(1,Y,S)\), choose
\begin{equation}
 K=\max\left(0,\left\lceil\log_2\frac{8M_0\sqrt n}{\eta}
                                          \right\rceil\right).
 \label{eq:panel-degree}
\end{equation}
On panel \(b\), the coefficients
\(\varrho(t_b)^ku^{(k)}(t_b)/k!\) have coordinate norm at most
\(M_0\sqrt n\), by Cauchy's formula on its disk.  Since
\(|t-t_b|/\varrho(t_b)\le1/2\), the degree-\(K\) tail is at
most \(M_0\sqrt n\,2^{-K}\le\eta/8\).  All coefficients span
one time-independent space of dimension at most \(J(K+1)\) per
curve.  Adjacent approximants need not agree at a panel boundary:
each approximates the same source there, and the comparison proof
uses source values, never derivatives of the piecewise approximant.

These coefficients can be computed from finite initialization jets,
including derivatives requested at later anchors.  Indeed \(D_1\)
contains a uniform rectangle about \([0,T_0]\) of radius
\(\varrho(0)>0\).  Use this radius in the explicit conformal map
\eqref{int:eq-0234-harmonic-initial-jet-map}, denoted here
\(t=\mathfrak t(\xi)\).  It satisfies \(\mathfrak t(0)=0\),
has nonzero derivative, maps the unit disk into that rectangle, and
maps \([0,\xi_*]\), \(\xi_*<1\), onto \([0,T_0]\).
For \(G(\xi)=u(\mathfrak t(\xi))\),
\[
 [\xi^j]G=\sum_{k=0}^j\frac{u^{(k)}(0)}{k!}
                              [\xi^j]\mathfrak t(\xi)^k.
\]
Every coefficient therefore uses finitely many zero-time derivatives.
They are obtained by differentiating the dense ODE at initialization;
in normalized form, if \(\dot\Theta=F(\Theta)\), the recurrence is
\(\Theta_{j+1}=(j+1)^{-1}[t^j]F(\sum_{k=0}^j\Theta_kt^k)\).
It uses initialized weights, training data/labels, passive inputs,
and the specified activation-derivative evaluator, not trained weights.

To see that later derivatives are also finite computations, choose
\(\xi_*<\rho_*<1\) and \(0<\delta_*<\rho_*-\xi_*\).
The degree-\(N\) Taylor sum \(G_N\) obeys
\[
 \sup_{|\xi|\le\rho_*}|G-G_N|
 \le\frac{M_0\sqrt n\,\rho_*^{N+1}}{1-\rho_*},\qquad
 |(G-G_N)^{(k)}(\xi_b)|
 \le\frac{k!M_0\sqrt n\,\rho_*^{N+1}}
             {\delta_*^k(1-\rho_*)}.
\]
The second inequality is Cauchy's derivative formula.  Applying
\((\mathfrak t'^{-1}\partial_\xi)^k\) converts these to physical
derivatives; its coefficients are finite at the finitely many anchors
for \(k\le K\).  A finite common \(N\) attains any prescribed
coefficient accuracy. In particular, coordinate tolerance
\(\eta/(128\sqrt n)\) for each scaled coefficient guarantees the
Euclidean tolerance used below. This proves finite initialization-only
provenance, not cheap arithmetic or well-conditioned numerical recovery.

Approximate each scaled base coefficient in Euclidean norm to
\(\eta/128\), and form its paired initialized image by the exact
initialized matrix action.  The base reconstruction error is at most
\((\eta/128)\sum_{k\ge0}2^{-k}=\eta/64\); initialized norm eight
makes its image error at most \(\eta/8\).  Adding Taylor tails keeps
both source errors below \(\eta\), while their image relation is
exact.  Thus no small coordinate error is incorrectly multiplied by
a width-dependent infinity norm.  An invertible conversion of each
local polynomial to Chebyshev coordinates changes neither its vector
span nor these paired actions; it is optional.

Add the constant vector, exact initialized training features and their
images, and all first-weight columns.  The resulting source dimension
per layer is at most
\begin{equation}
 R=2m+d+1+2(2m+p)J(K+1).
 \label{eq:panel-source-rank}
\end{equation}
The selection proof \S\ref{int:harmonic-selection-proof} supplies
\(q_j\le9R\), exact normalized source isometry, and metrics
\(D_j/4\preceq M_j\preceq D_j\),
\(\mathbf1^\top D_j\mathbf1\le4\).
The same selected initialized mixer and its metric adjoint preserve the
forward and reverse image pairs.  Exact training additions preserve the
initial training Gram.  These verify every initialization hypothesis of
\S\ref{int:harmonic-runtime-definition}.

\subsection{Autonomous runtime, all-time error, and retained inventory}
\label{sec:panel-runtime-proof}

Use the moving arrays and corrected readout
\eqref{int:eq-0061-harmonic-state}--\eqref{int:eq-0064-harmonic-autonomous-optimizer},
with training indices only in the deficit and Gram updates.  The
training prediction is exactly \(y-c_C\); the passive inputs require
neither deficits nor a Gram inverse.  The fitting proof
\S\ref{int:harmonic-fitting-proof} applies because the metric bounds,
initial operator bounds, exact Gram, and label allowance have just
been verified.  In particular its normalized training Gram remains
above \(\lambda/4\), residuals decay, and parameters converge.

For clarity, the comparison does not follow from fitting alone.
Source coordinate error \(\eta\) has dense RMS at most \(\eta\)
and selected metric norm at most \(2\eta\).  Exact isometry therefore
gives, for source RMS bounds \(A_1,A_2\), pairing error at most
\(3\eta(A_1+A_2)+9\eta^2\), and paired initialized-action error
at most \(18\eta\).  Learned actions integrate those defects
against residuals.  The source-energy proof
\S\ref{int:harmonic-source-energy-proof} controls the selected readout
by \(O(Y/\sqrt\lambda)\).  With
\(V_C=[h_{C,a}^{(L)}/\sqrt m]_a\),
\(Q_C=V_C^*V_C\), \(T_C=V_CQ_C^{-1}\), and
\(e=(c_C-c_n)/\sqrt m\), the error variable
\(w_C-R_Lw_n+T_Ce\) has the exact cancellation
\(+2V_Ce-2T_CQ_Ce=0\).
This is \eqref{int:eq-0252-harmonic-exact-readout-cancellation}, not
the false assertion that \(T_CV_C^*\) is the identity on all features.
Remaining terms are handled by signed damping and residual-weighted
integrals as in \S\ref{int:harmonic-cancellation-proof}.
Backward terms and Gram inverses involve training points only;
passive forward subtraction uses the declared source curves and a
maximum over them introduces no factor \(p\).

Consequently the identical recurrence certificate
\eqref{int:eq-0069-harmonic-horizon-and-tolerance-certificate} gives
\[
 \sup_{t\in[0,\infty]}\max_{i\le m+p}
       |f_C(t,x_i)-f_n(t,x_i)|
 \le\mathcal A_n\eta+\mathcal D e^{-8\ell}.
\]
Here \(\mathcal A_n\) is exactly
\eqref{int:eq-0068-harmonic-exact-error-coefficient}, with coefficient
\(64\) in its \(\sqrt\ell\) exponent, and \(\mathcal D\) is
\eqref{int:eq-0057-harmonic-tail-coefficient}.  The entire label
interval permits the completely numerical substitutes
\begin{align*}
 \overline{\mathcal A}_n
 &=2000e^{44}\beta^{42L}\frac Y\lambda
            (1+\lambda^{-1/2})(1+\sqrt\ell)e^{64\sqrt\ell},\\
 \overline{\mathcal D}
 &=236\beta^{9L}\frac Y\lambda(1+\lambda^{-1/2}),
\end{align*}
by \eqref{int:eq-0264-harmonic-numerical-prefactor-bound}--\eqref{int:eq-0265-harmonic-numerical-tail-bound}.
For the stronger target \(Y/n\), take
\begin{equation}
 \eta=\min\{1,Y,S,Y/(2n\overline{\mathcal A}_n)\},\qquad
 \overline{\mathcal D}e^{-8\ell}\le Y/(2n).
 \label{eq:panel-tolerance}
\end{equation}
Source comparison is used only through \(T_0\).  Afterward each flow
is compared to its own value at \(T_0\), using the separate fitted
tails \eqref{int:eq-0240-harmonic-endpoint-tail}.  Both flows continue
to evolve; neither is frozen or played back.  Their limits give the
same endpoint bound.

The complete retained real-coordinate inventory is at most
\begin{equation}
 1020(L+1)R^2+36R+10m(d+1)+pd+(m+p)+D_{\rm alg}.
 \label{eq:panel-exact-inventory}
\end{equation}
Indeed \eqref{int:eq-0082-harmonic-all-retained-inventory} counts
moving matrices, initialized copies, metrics and inverse caches,
training/Gram buffers, readout workspace, and training data.
Four reusable neuron buffers add at most \(36R\); passive input and
panel output storage add \(pd+(m+p)\).  The separately charged
\(D_{\rm alg}\) contains the scalar evaluator, program, and evaluator
workspace; it may not conceal width-dependent source arrays.
Dense weights, coefficient vectors, jets, panel anchors, and selected
index lists are discarded.  There is no retained temporal forcing or
panel schedule.  The count concerns real coordinates, not bits,
conditioning, setup work, or numerical integration step count.

An exact input-span reduction improves the dimension overhead.  Let
\(U_0\in\mathbb R^{d\times k}\) have orthonormal columns spanning
all declared inputs, \(k\le\min(d,m+p)\).  Replace inputs by
\(U_0^\top v_i\) and the first matrix by \(AU_0\).
Every first-layer update lies in this span, so
\(A(t)(I-U_0U_0^\top)=A(0)(I-U_0U_0^\top)\).
The reduced and original coupled trajectories agree on the panel by
induction through the layers.  Each Gaussian row of \(A(0)U_0\) has
covariance \(I_k\), and all panel inner products and the population
Gram are unchanged.  Use
\(R=2m+k+1+2(2m+p)J(K+1)\) and replace \(d\) by \(k\) in
\eqref{eq:panel-exact-inventory}; a safe additional charge for the
map and both original/transformed panel data is \(3(m+p)d\).
This is exact on the declared panel, not a whole-sphere reduction.

Finally impose the directly checkable eventual degree gate
\begin{equation}
 K+1\le8\log(en).
 \label{eq:panel-degree-gate}
\end{equation}
It follows eventually from \eqref{eq:panel-degree} and
\eqref{eq:panel-tolerance}.  Squaring the three terms of
\eqref{eq:panel-panel-count} with
\((u+v+w)^2\le3(u^2+v^2+w^2)\), and using the exact input-span
count, gives a sufficient envelope
\begin{equation}
 \begin{split}
 C\beta^{CL}(m+p)^2\bigg[
 &\left(\frac{Ym}{\gamma}\right)^4[\log(en)]^3\\
 &+\left(\frac m\gamma\right)^2[\log(en)]^2
                   [\log(e+\log(en))]^2\bigg]
 +C(m+p)d+D_{\rm alg}.
 \end{split}
 \label{eq:panel-final-inventory}
\end{equation}
Here only the final envelope uses an absolute numerical \(C\).
To verify its parameter dependence directly, the recurrences above
give \(H_j\le\beta^{3j}\),
\(\tau_j\le\beta^{5L-2j+1}\),
\(U/a\le\beta^{60L}\), and \(\mathcal K\le\beta^{14L}\).
For example the intermediate exponent bounds
\(A_*\le\beta^{8L}\), \(H_*\le\beta^{11L}\),
\(D_0\le\beta^{24L}\), \(K_{\rm src}\le\beta^{32L}\),
\(\max e_j\le\beta^{17L}\), \(T_Q\le\beta^{22L}\),
and \(U\le\beta^{57L}\) follow by geometric unrolling, using
\(10s\le\beta^2\) and \(L\le\beta^L\).
Also \(\mathcal K/\lambda\ge1\), so the lower-order initial ranks
and \(\ell^2\) term are absorbed in the second displayed term.
The bound \(S\le1\) is used only for these source coefficients;
the explicit factor \((Ym/\gamma)^4\) has not been removed by
spending the label allowance.  The two sample/gap factors in
\eqref{eq:panel-final-inventory} must not be identified or compared
by treating \(m\) and \(1/\gamma\) as equal.

\subsection{Assembly of the headline theorem}
\label{sec:panel-headline-proof}

For \(m\ge2\), the dense-variability proof in
\S\ref{int:dense-innovation-lower} and
\S\ref{int:dense-trajectory-lower} selects a deterministic
\emph{training index}, not an arbitrary sphere point.  To recall the
reason, for one initialized feature vector \(H\in\mathbb R^m\),
the positive Gram gap and fourth-moment argument give
\[
 \operatorname{tr}\operatorname{Cov}(H(y^\top H))
 \ge\frac{\gamma^3\mu_{2,L}}{16\mu_4}\|y\|_2^2.
\]
At least one training coordinate carries this variance divided by
\(m\).  The initialized Gram CLT and the independent second run
then give an onset-velocity fluctuation there; the finite-query
time-analytic argument in \S\ref{int:dense-trajectory-lower}
converts it into an actual positive-time discrepancy.  Thus for every
fixed \(0<\delta_0<1\), eventually
\[
 \Pr\left\{\|f_n-\widetilde f_n\|_{\mathcal X}\ge
       c_{\phi,L,\delta_0}\frac{Y\sqrt\gamma}
                    {\sqrt n[\log(en)]^{5/2}}\right\}
       \ge1-\delta_0
\]
for the declared panel \(\mathcal X\), with exactly the positive
coefficient in \S\ref{int:dense-trajectory-lower}.  This is not an
endpoint lower bound.  In particular, taking \(\delta_0=1/2\)
puts the displayed lower scale below the fixed \(99.99\%\) quantile
\(b_n(\mathcal X)\): every smaller threshold has cumulative
probability at most \(1/2\).  The error \(Y/n\) proved above is
eventually below that scale and hence below \(3b_n(\mathcal X)\).
More strongly its ratio to realized independent-dense panel discrepancy
tends to zero in probability, by first fixing arbitrary \(\delta_0\)
and then letting \(n\) grow.  Independence between the compressed
construction and the lower event is unnecessary; a union bound suffices.

This proves the Logarithmic row and its all-time comparison in
Theorem~\ref{thm:compression}.  The other two rows use their unchanged
proofs: Legendre's signed-stability and optimized-order arguments are
in \S\ref{int:legendre-stability-proof} and
\S\ref{int:legendre-inverse-proof}; Harmonic's joint time/input
source count and inverse are in
\S\ref{int:harmonic-expansion-proof} and
\S\ref{int:harmonic-budget-inversion-proof}, using the label-explicit
inverse \eqref{int:eq-0021-harmonic-inverse} rather than its optional
small-label simplification.  Their full-sphere errors also become
negligible relative to the corresponding dense-pair discrepancy at
the displayed orders.  Fixed-confidence intersections cost only a
change in the qualitative sufficient width, not in the storage powers.
The query domains remain distinct: a declared-panel certificate is not
an unseen-input theorem.  No input-span condition, passive-label use,
bounded-activation-value assumption, or additional label cap was added.

For \(m=1\) the additive \(Y/n\) theorem still holds, but positive
uncentered Gram alone need not give nonzero dense variability (a
constant final activation is a counterexample).  The multiplicative
comparison is therefore stated for \(m\ge2\).  The strict
initialization-only compiler above may be expensive; neither its
retained-coordinate theorem nor its finite-jet provenance asserts
near-linear or near-quadratic setup for this new panel construction.
